\documentclass[11pt,a4paper]{article}
\usepackage[margin=23mm]{geometry}
\usepackage{amsmath,amssymb,booktabs,tabularx,array}
\usepackage{graphicx}
\usepackage{enumitem,microtype}
\usepackage[hidelinks]{hyperref}
\setlength{\parindent}{0pt}
\setlength{\parskip}{4pt}
\setlist[itemize]{leftmargin=1.4em,itemsep=3pt,topsep=4pt}
\setlist[enumerate]{leftmargin=1.6em,itemsep=4pt,topsep=4pt}
\renewcommand{\arraystretch}{1.18}
\newcommand{\E}{\mathbb E}
\newcommand{\ind}{\mathbf 1}
\newcommand{\R}{\mathbb R}
\newcommand{\lead}[1]{\medskip\textbf{#1}\par}
\newcommand{\takeaway}[1]{\medskip\noindent\fbox{\parbox{\dimexpr\linewidth-2\fboxsep-2\fboxrule\relax}{#1}}\medskip}
\title{Cooperative Target Assignment}
\author{}
\date{}
\begin{document}
\maketitle
\vspace{-16mm}

\takeaway{\textbf{Goal:} Assign five drones to five targets to maximize the total target value removed within 100 seconds. Each completed strike consumes one drone.}

\section{Mission setup}
\begin{center}
\begin{tabularx}{\linewidth}{@{}lX@{}}
\toprule
Quantity & Setting \\
\midrule
Drones / targets & $N=5$ / $M=5$ \\
Map / coordinate unit & $10\times10$ km / 100 m \\
Horizon / time step & $H=100$ seconds / $\Delta t=1$ second \\
Drone speed & $v_{\max}=4/9$ units/s $=160$ km/h \\
Target motion & Stationary \\
Strike range / duration & 5 m / $K=10$ seconds per drone \\
Strike success probability & 1 \\
Available information & All target positions, types, and remaining lives; friendly states and previous intentions \\
Decision & One target per active drone, every second during approach \\
\bottomrule
\end{tabularx}
\end{center}

\lead{Target types}
\begin{center}
\begin{tabular}{@{}ccccc@{}}
\toprule
Type & Count & Initial lives & Remaining value & Strike capacity \\
\midrule
1 & 2 & 2 & $5\to2.5\to0$ & Up to 2 drones \\
2 & 2 & 1 & $2\to0$ & 1 drone \\
3 & 1 & 1 & $1\to0$ & 1 drone \\
\bottomrule
\end{tabular}
\end{center}

\lead{Training layouts}
\begin{itemize}
\item Friendly starting cluster: radius 100 m.
\item Two target formations: three targets and two targets; radius 650 m each.
\item First formation centre: 2.5--5.8 km from the friendly centre.
\item Second centre: 55\% along that route, offset laterally by 800 m.
\end{itemize}

The success threshold $B$ is the first formation's initial value.
Training samples each row below with equal probability.
\begin{center}
\begin{tabular}{@{}ccc@{}}
\toprule
$B$ & Formation 1 types & Formation 2 types \\
\midrule
5 & $(2,2,3)$ & $(1,1)$ \\
8 & $(1,2,3)$ & $(1,2)$ \\
9 & $(1,2,2)$ & $(1,3)$ \\
11 & $(1,1,3)$ & $(2,2)$ \\
12 & $(1,1,2)$ & $(2,3)$ \\
\bottomrule
\end{tabular}
\end{center}

\newpage
\section{State space}
\takeaway{The state collects remaining time, physical states, and the team's previous intentions. The critic receives \textbf{111 scalar entries}.}

Indices: drone $i,m\in\mathcal I=\{1,\ldots,N\}$;
target $k\in\mathcal K=\{1,\ldots,M\}$; time $t\in\{0,\ldots,H\}$.

\begin{center}
\begin{tabularx}{\linewidth}{@{}lXl@{}}
\toprule
Symbol & Meaning & Domain \\
\midrule
$r_t^i$ & Drone position & $[0,L_{\rm map}]^2$ \\
$\sigma_t^i$ & Active-drone indicator & $\{0,1\}$ \\
$\zeta_t^i$ & Elapsed strike steps; zero during approach & $\{0,\ldots,K-1\}$ \\
$v_t^i$ & Measured world-frame velocity & $\R^2$ \\
$q^k$ & Stationary target position & $[0,L_{\rm map}]^2$ \\
$c^k$ & One-hot target type & $\{e_1,e_2,e_3\}$ \\
$L_k(t)$ & Remaining lives & $\{0,\ldots,L_k(0)\}$ \\
$P_{t-1}^{ik}$ & Previous probability of selecting target $k$ & $[0,1]$ \\
$C_{t-1}^{ik}$ & Previous actual selection of target $k$ & $\{0,1\}$ \\
\bottomrule
\end{tabularx}
\end{center}

\lead{Normalization}
With $L_{\rm map}=100$ and $K=10$,
\begin{align*}
\tau_t&=(H-t)/H,& \bar r_t^i&=r_t^i/L_{\rm map},& \bar q^k&=q^k/L_{\rm map},\\
\bar v_t^i&=v_t^i/v_{\max},& \bar\zeta_t^i&=\zeta_t^i/K,& \bar L_k(t)&=L_k(t)/2.
\end{align*}

\lead{Global state}
\[
s_t=\bigl(\tau_t,\{d_t^i\}_{i=1}^N,\{y_t^k\}_{k=1}^M,P_{t-1},C_{t-1}\bigr),
\]
\[
\underbrace{d_t^i=(\bar r_t^i,\sigma_t^i,\bar\zeta_t^i,\bar v_t^i)}_{\text{drone record: 6 entries}},
\qquad
\underbrace{y_t^k=(\bar q^k,c^k,\bar L_k(t))}_{\text{target record: 6 entries}}.
\]
\begin{center}
\begin{tabular}{@{}lrrrrr@{}}
\toprule
Block & Time & Drones & Targets & $P$ & $C$ \\
\midrule
Entries & 1 & $6N=30$ & $6M=30$ & $NM=25$ & $NM=25$ \\
\bottomrule
\end{tabular}
\end{center}
\[
\boxed{\mathcal S\subset\R^{1+6N+6M+2NM}=\R^{111}}
\]

\lead{Intention and strike participation}
\begin{itemize}
\item Initially: $P_{-1}^{ik}=1/M$ and $C_{-1}^{ik}=0$.
\item Each active drone subsequently has one probability row in $P$ and one selected target in $C$.
\item Inactive drones: zero intention rows and velocity; masked drone nodes.
\end{itemize}
\[
S_t^{ik}=\ind\{\sigma_t^i=1,\ \zeta_t^i>0\}\,C_{t-1}^{ik}
\quad\text{indicates an ongoing strike on target }k.
\]

\newpage
\section{Actions, transitions, and objective}
\lead{Action space}
\[
a_{i,t}\in\mathcal K,\qquad
\boldsymbol a_t\in\mathcal K^N,\qquad |\mathcal K^N|=5^5=3125.
\]
For target $k$, let $b_k=2$ for Type 1 and $b_k=1$ otherwise.
\[
\kappa_k(t)=\min\{b_k,L_k(t)\},\qquad n_k(t)=\sum_i S_t^{ik}.
\]
\begin{center}
\begin{tabularx}{\linewidth}{@{}lX@{}}
\toprule
Drone status & Target selection \\
\midrule
Approaching & $\mathcal A_i(s_t)=\{k:L_k(t)>0,\ n_k(t)<\kappa_k(t)\}$ \\
Striking & Keep the current target until completion \\
Inactive & Action has no physical effect \\
Empty available set & Use target index 1 as a categorical fallback \\
\bottomrule
\end{tabularx}
\end{center}

\lead{State transition}
\begin{enumerate}
\item \textbf{Move} each active drone toward its selected target.
\item \textbf{Admit arrivals} within 5 m into free strike slots. Existing participants retain their slots; excess simultaneous arrivals are selected randomly.
\item \textbf{Advance strikes} by one second per participant, including new arrivals.
\item \textbf{Complete strikes}: after ten seconds, remove one target life per successful drone and consume the completing drones.
\item \textbf{Update} reward, velocity, remaining time, and previous intentions.
\end{enumerate}
\[
r_{t+1}^i=r_t^i+
\min\!\left(1,\frac{v_{\max}\Delta t}{\|q^{a_{i,t}}-r_t^i\|}\right)
(q^{a_{i,t}}-r_t^i).
\]
The displacement is zero at the destination. Target positions remain fixed.

\lead{Reward and mission success}
Let $V_k(t)$ denote remaining target value, as listed in the target-type table.
\begin{align*}
D_t&=\sum_k[V_k(0)-V_k(t)],& r_t&=D_{t+1}-D_t,\\
B&=\sum_{k\in\text{formation 1}}V_k(0),& \text{success}&=\ind\{D_H\ge B\}.
\end{align*}
\takeaway{The team reward is the decrease in remaining target value. Both formations contribute to the score.}

\lead{Finite-horizon objective}
\[
\boxed{J(\theta)=\E_{\pi_\theta}\!\left[\sum_{t=0}^{H-1}r_t\right]=\E_{\pi_\theta}[D_H]},
\qquad \gamma=1.
\]
Episodes end at $H=100$. For the remaining mission,
\[
V^\pi(s_t)=\E_\pi\!\left[\sum_{u=t}^{H-1}r_u\mid s_t\right],\qquad
Q^\pi(s_t,\boldsymbol a_t)=\E_\pi\!\left[\sum_{u=t}^{H-1}r_u\mid s_t,\boldsymbol a_t\right].
\]

\newpage
\section{Observation and policy}
\takeaway{The actor receives relative geometry and shared previous intentions, then assigns a selection probability to each available target.}

\lead{Local observation}
\[
x_t^{ik}=\left(\frac{q^k-r_t^i}{L_{\rm map}},c^k,\bar L_k(t)\right)\in\R^6,
\qquad \delta_t^{im}=\frac{r_t^m-r_t^i}{L_{\rm map}}\in\R^2.
\]
\[
\rho_t^i=\sigma_t^i\ind\{\zeta_t^i=0\}
\quad\text{indicates an active drone available to choose a target.}
\]
\begin{center}
\begin{tabularx}{\linewidth}{@{}lX@{}}
\toprule
Observation block & Contents \\
\midrule
Time and self & $\tau_t,\sigma_t^i,\rho_t^i,\bar v_t^i$ \\
Each target $k$ & Relative position, one-hot type, remaining lives: $x_t^{ik}$ \\
Each teammate $m\ne i$ & Relative position $\delta_t^{im}$, activity $\sigma_t^m$, availability $\rho_t^m$, velocity $\bar v_t^m$ \\
Shared intentions & Previous probabilities $P_{t-1}$ and actual choices $C_{t-1}$ \\
\bottomrule
\end{tabularx}
\end{center}
\[
\begin{split}
o_t^i=\bigl(&\tau_t,\sigma_t^i,\rho_t^i,\bar v_t^i,\{x_t^{ik}\}_k,\\
&\{(\delta_t^{im},\sigma_t^m,\rho_t^m,\bar v_t^m)\}_{m\ne i},P_{t-1},C_{t-1}\bigr).
\end{split}
\]

\lead{Target scores and probabilities}
The shared actor assigns a score $\ell_t^{ik}$ to each candidate target:
\[
\ell_t^{ik}=f_\theta(o_t^i,k),\qquad
\pi_{\theta,i}(k\mid o_t^i)=
\frac{\exp(\ell_t^{ik})}{\sum_{j\in\mathcal A_i(s_t)}\exp(\ell_t^{ij})}.
\]
\begin{center}
\begin{tabularx}{\linewidth}{@{}lX@{}}
\toprule
Quantity & Role \\
\midrule
$\theta$ & Actor parameters shared across drones and targets \\
$\ell_t^{ik}$ & Learned preference for assigning drone $i$ to target $k$ \\
$\mathcal A_i(s_t)$ & Available targets after applying strike-capacity constraints \\
$\pi_{\theta,i}$ & Target-selection probabilities that sum to one over available targets \\
\bottomrule
\end{tabularx}
\end{center}

\lead{Action selection}
\begin{itemize}
\item Training: each drone samples from its distribution.
\item Evaluation: each drone selects its highest-scoring available target; the score resolver handles assignment conflicts.
\end{itemize}
The joint training policy factorizes as
\[
\pi_\theta(\boldsymbol a_t\mid\boldsymbol o_t)
=\prod_{i=1}^{N}\pi_{\theta,i}(a_{i,t}\mid o_t^i).
\]

\newpage
\section{Actor--critic learning}
\takeaway{The critic learns the return of a joint assignment. The actor learns which individual choices improve that return.}

\begin{center}
\begin{tabularx}{\linewidth}{@{}lXX@{}}
\toprule
 & Actor & Critic \\
\midrule
Input & Drone observation $o_t^i$ & Global state and joint choice $(s_t,\boldsymbol a_t)$ \\
Output & Target probabilities $\pi_{\theta,i}$ & Predicted remaining team return $Q_\phi$ \\
Learning signal & Counterfactual advantage & Collected rewards and a continuation estimate \\
Parameters & $\theta$ & $\phi$; slowly updated target copy $\bar\phi$ \\
\bottomrule
\end{tabularx}
\end{center}
\[
\boxed{\text{Collect}}\ \longrightarrow\
\boxed{\text{Fit critic}}\ \longrightarrow\
\boxed{\text{Assign credit}}\ \longrightarrow\
\boxed{\text{Update actor}}
\]

\subsection{Critic objective}
\begin{enumerate}
\item Collect states, observations, joint actions, action probabilities, rewards, and episode endings using $\pi_{\mathrm{old}}$.
\item Compute a return target $G_t$ backward through the rollout.
\item Fit $Q_\phi(s_t,\boldsymbol a_t)$ to $G_t$ and update the target copy.
\end{enumerate}
\[
G_t=r_t+\gamma(1-d_t)
\bigl[(1-\lambda)Q_{\bar\phi}(s_{t+1},\boldsymbol a_{t+1})+\lambda G_{t+1}\bigr].
\]
\begin{center}
\begin{tabularx}{\linewidth}{@{}lX@{}}
\toprule
Term & Meaning \\
\midrule
$r_t$ & Team reward observed immediately after the joint action \\
$Q_{\bar\phi}(s_{t+1},\boldsymbol a_{t+1})$ & Target critic's estimate of the remaining return \\
$G_{t+1}$ & Return target carried backward from the next step \\
$\lambda=0.97$ & Weight controlling how far collected rewards propagate backward \\
$d_t$ & 1 at the mission horizon, which sets continuation to zero \\
\bottomrule
\end{tabularx}
\end{center}
At a rollout cutoff, the target critic supplies the continuation value; $\gamma=1$.
\[
\mathcal L_Q=\E\bigl[\operatorname{Huber}_{10}(Q_\phi(s_t,\boldsymbol a_t)-G_t)\bigr],
\qquad \bar\phi\leftarrow0.99\bar\phi+0.01\phi.
\]
Huber loss penalizes prediction error quadratically near zero and linearly for large errors.
The target copy is updated after each critic optimizer step.

\subsection{Counterfactual baseline}
\begin{enumerate}
\item Hold the state and all other drones' choices $\boldsymbol a_{-i}$ fixed.
\item Evaluate each alternative target $u$ for drone $i$ with the critic.
\item Average those values using the collected action probabilities.
\item Subtract this average from the value of the collected joint action.
\end{enumerate}
\begin{align*}
b_i(s,\boldsymbol a_{-i})
&=\sum_{u\in\mathcal K}\pi_{\mathrm{old},i}(u\mid o_i)
Q_\phi\bigl(s,(\boldsymbol a_{-i},u)\bigr),\\
A_i(s,\boldsymbol a)&=Q_\phi(s,\boldsymbol a)-b_i(s,\boldsymbol a_{-i}).
\end{align*}
$A_i$ measures how the selected target compares with
that drone's probability-weighted alternatives, given the other drones' choices.

\newpage
\subsection{PPO objective}
Normalize advantages over valid actor decisions and keep them fixed during the actor update:
\[
\widehat A_i=\frac{A_i-\operatorname{mean}(A)}{\operatorname{std}(A)+10^{-5}},
\qquad
\omega_i=\frac{\pi_{\theta,i}(a_i\mid o_i)}{\pi_{\mathrm{old},i}(a_i\mid o_i)}.
\]
\begin{align*}
\mathcal L_\pi
&=-\E\!\left[\min\!\left(\omega_i\widehat A_i,
\operatorname{clip}(\omega_i,1-\epsilon,1+\epsilon)\widehat A_i\right)\right]\\
&\quad-\beta\E[\mathcal H(\pi_i)].
\end{align*}
\begin{center}
\begin{tabularx}{\linewidth}{@{}lX@{}}
\toprule
Component & Effect on learning \\
\midrule
Positive $\widehat A_i$ & Encourages the collected choice \\
Negative $\widehat A_i$ & Discourages the collected choice \\
Probability ratio $\omega_i$ & Measures the change from the rollout policy \\
Clipping with $\epsilon=0.2$ & Caps the surrogate's incentive for a large change in probability \\
Entropy bonus, $\beta=0.01$ & Encourages continued exploration across available targets \\
Decision mask & Includes the steps at which a drone actually chooses a target \\
\bottomrule
\end{tabularx}
\end{center}

\lead{Update schedule}
\begin{enumerate}
\item Keep the collected return targets fixed while fitting the critic.
\item Use the fitted critic to compute counterfactual advantages.
\item Keep those advantages fixed while updating the actor.
\item Collect a new rollout with the updated actor.
\end{enumerate}
The actor and critic use separate optimizers and separate parameters.

\lead{Training parameters}
\begin{center}
\begin{tabularx}{\linewidth}{@{}lX@{}}
\toprule
Parameter & Value \\
\midrule
Environments / rollout length & 32 / 200 steps per environment \\
Actor / critic epochs & 10 / 40 \\
Actor / critic minibatch & 256 drone rows / 256 joint transitions \\
Optimizer & Adam; learning rate $5\times10^{-4}$, epsilon $10^{-5}$ \\
Discount / trace coefficient & $\gamma=1$ / $\lambda=0.97$ \\
Gradient norm limit & 10 \\
Training seed & 7 \\
\bottomrule
\end{tabularx}
\end{center}
The training counter increases once per active drone per simulation step.

\newpage
\section{Shared networks for variable drone and target counts}
\takeaway{Represent each entity with a fixed-width record, apply the same learned functions to every record, and combine records by weighted sums. Changing $N$ or $M$ changes how many records are processed.}

\subsection{Encoding and attention}
\begin{center}
\begin{tabularx}{\linewidth}{@{}lX@{}}
\toprule
Term & Meaning in this model \\
\midrule
Record & A fixed list of measured attributes for one entity or relation \\
Encoding & A learned function that converts a record into 64 numerical features \\
Token & One encoded record, represented by that feature vector \\
Attention & Computing an input-dependent weight for each source record and summing its weighted contribution \\
Mask & A zero/one indicator that removes inactive or absent entries from a sum \\
\bottomrule
\end{tabularx}
\end{center}

Let $h$ represent the recipient and $z_j$ a source record, both after encoding.
Each attention operation performs the following steps:
\begin{enumerate}
\item \textbf{Compare}: transform $h$ and $z_j$ into vectors used to measure their relevance.
\item \textbf{Weight}: convert the comparison into a number $g_j$ between zero and one.
\item \textbf{Summarize}: multiply the source's learned contribution $b_j$ by $g_j$, then sum over valid sources.
\item \textbf{Update}: combine this summary with the recipient's current features.
\end{enumerate}

For one parallel comparison channel,
\begin{align*}
q&=W_qh, & k_j&=W_kz_j,\\
g_j&=\operatorname{sigmoid}\!\left(\frac{q^\top k_j}{\sqrt{d_h}}\right),
& b_j&=\operatorname{softplus}(W_vz_j),\\
m&=\sum_j\mu_j g_j b_j,
& \operatorname{sigmoid}(x)&=\frac{1}{1+e^{-x}}.
\end{align*}
Here $\mu_j\in\{0,1\}$ is the validity mask and
$\operatorname{softplus}(x)=\log(1+e^x)$ is applied to each coordinate.
Bias terms are included in the learned transformations.

\begin{itemize}
\item Four parallel channels, called \textbf{heads}, each produce a 16-entry summary.
\item Joining the four summaries gives $m_{\rm all}\in\R^{64}$.
\item A small learned function $F$ updates the recipient:
$h' = h+F([h,m_{\rm all}])$.
\item $[h,m_{\rm all}]$ means joining the two vectors end to end.
\end{itemize}
Each weighted sum returns 16 entries per head,
regardless of the number of sources. Summing the contributions lets the
aggregate reflect how many entities contribute.

\newpage
\subsection{Shared target-scoring actor}
\lead{Fixed-width inputs}
For acting drone $i$ and candidate target $k$,
\begin{align*}
u_t^{ik}&=(x_t^{ik},\tau_t,P_{t-1}^{ik},C_{t-1}^{ik},\sigma_t^i,\rho_t^i,\bar v_t^i)\in\R^{13},\\
w_t^{imk}&=(\delta_t^{im},\rho_t^m,\bar v_t^m,P_{t-1}^{mk},C_{t-1}^{mk})\in\R^7.
\end{align*}
\begin{center}
\begin{tabularx}{\linewidth}{@{}lcX@{}}
\toprule
Record & Width & Contents \\
\midrule
Candidate $u_t^{ik}$ & 13 & Target attributes, time, own previous intention, own activity and velocity \\
Teammate $w_t^{imk}$ & 7 & Relative position, availability, velocity, previous intention for target $k$ \\
\bottomrule
\end{tabularx}
\end{center}

\lead{Target scoring}
\begin{enumerate}
\item Encode the candidate record and each teammate record into 64 features.
\item Use the candidate as the recipient and active teammates as sources.
\item Apply two successive attention updates to the candidate features.
\item Convert the updated candidate features into one scalar score $\ell_t^{ik}$.
\end{enumerate}
\[
\boxed{u_t^{ik}\in\R^{13}}\ \longrightarrow\
\boxed{h^{ik}\in\R^{64}}\ \xrightarrow{\text{read teammates twice}}\
\boxed{\ell_t^{ik}\in\R}.
\]
The candidate representation is the \textbf{query}: it determines which
teammate information matters for scoring that target.

\lead{Variable input cardinality}
\begin{center}
\begin{tabularx}{\linewidth}{@{}lX@{}}
\toprule
Change & Effect on the computation \\
\midrule
More drones & More teammate records enter each weighted sum \\
More targets & Repeat the same scoring function for more candidate records \\
Target selection & Apply the masked softmax over the resulting $M$ scores \\
Trainable functions & The candidate encoder, teammate encoder, attention updates, and scalar output map retain the same dimensions \\
\bottomrule
\end{tabularx}
\end{center}

\lead{Parameter sharing}
\[
E_{\rm candidate}:\R^{13}\to\R^{64},\qquad
E_{\rm peer}:\R^7\to\R^{64},\qquad
f_{\rm score}:\R^{64}\to\R.
\]
Each function is reused across all drones and targets. Two attention layers
have their own learned parameters, each shared across all candidate computations.

\newpage
\subsection{Graph-based team critic}
\lead{Fixed-width inputs}
\begin{center}
\begin{tabularx}{\linewidth}{@{}lccX@{}}
\toprule
Record & Number & Width & Contents \\
\midrule
Drone $d_t^i$ & $N$ & 6 & Position, activity, strike progress, velocity \\
Target $y_t^k$ & $M$ & 6 & Position, type, remaining lives \\
Relation $e_t^{ik}$ & $NM$ & 5 & Current choice, previous intention, relative position \\
\bottomrule
\end{tabularx}
\end{center}
\[
e_t^{ik}=\bigl(\ind\{a_{i,t}=k\},P_{t-1}^{ik},C_{t-1}^{ik},\bar q^k-\bar r_t^i\bigr).
\]
An entity record is called a \textbf{node}; a drone--target relation is called
an \textbf{edge}. The collection forms the critic's graph.

\lead{Relation-conditioned attention}
A recipient uses the same weighted-sum operation, now including its relation to each source:
\[
k_j=W_kz_j+W_{ek}e_j,\qquad
b_j=\operatorname{softplus}(W_vz_j+W_{ev}e_j).
\]
The relation influences both the relevance weight and the information sent.
The current assignment enters through $\ind\{a_{i,t}=k\}$.

\lead{Message passing and team readout}
\begin{enumerate}
\item Encode each drone and target into 64 features, adding a learned representation of remaining time.
\item Each target combines information from active drones using their relations.
\item Each drone combines information from the updated targets. Relative positions reverse direction for this exchange.
\item Repeat steps 2--3 once more.
\item Combine all valid drone and target features into one 64-entry team summary, using weights conditioned on remaining time.
\item Join that summary with the 64-entry time representation and map the resulting 128 entries to one scalar $Q_\phi$.
\end{enumerate}
\[
\boxed{N+M\text{ entity records}}\ \longrightarrow\
\boxed{\text{Team summary}\in\R^{64}}\ \longrightarrow\
\boxed{Q_\phi\in\R}.
\]

\lead{Parameter dimensions and permutation symmetry}
\begin{center}
\begin{tabularx}{\linewidth}{@{}lX@{}}
\toprule
Property & Reason \\
\midrule
Same parameter dimensions for any $N,M$ & Encoders process one fixed-width record at a time; weighted sums have a fixed output width \\
Actor scores follow entity reordering & The same scorer is applied to every candidate; teammate information is summed \\
Team value is invariant to entity reordering & Relations follow their entities, and the final team summary is a sum over valid entities \\
\bottomrule
\end{tabularx}
\end{center}
Let $\mathcal R$ consistently relabel drone and target records, intentions,
relations, and actions, with index permutations $\pi_D$ and $\pi_T$.
\[
\ell^{\pi_D(i),\pi_T(k)}(\mathcal R s)=\ell^{ik}(s),\qquad
Q_\phi(\mathcal R s,\mathcal R\boldsymbol a)=Q_\phi(s,\boldsymbol a).
\]
Here $\ell^{ik}(s)$ denotes the actor score formed from drone $i$'s observation
of $s$. The actor is permutation equivariant; the scalar critic is permutation invariant.

\newpage
\section{Evaluation}
\begin{itemize}
\item Scenarios are sampled from the training distribution, using the same
formation-layout and target-type sampling rules.
\item All policies and checkpoints are evaluated on the same 100 held-out
scenarios, with seeds 30,000--30,099.
\item Checkpoints are evaluated every 100,000 active-drone transitions.
\end{itemize}

\lead{Metrics}
\begin{center}
\begin{tabularx}{\linewidth}{@{}lX@{}}
\toprule
Metric & Definition \\
\midrule
Final score & $D_H$ \\
Team return & $\sum_{t=0}^{H-1}r_t=D_H$ \\
Mission success rate & Episode average of $\ind\{D_H\ge B\}$ \\
Destroyed fraction & $M^{-1}\sum_k\ind\{L_k(H)=0\}$ \\
Score AUC & $H^{-1}\sum_{t=1}^{H}D_t/B$; rewards earlier accumulation of score \\
\bottomrule
\end{tabularx}
\end{center}

\lead{Baseline allocation}
Nearest-first and type-priority use the following procedure:
\begin{enumerate}
\item Retain valid approach assignments and reserve their target slots.
\item For each candidate target, form a team from the closest free drones to fill its unreserved slots.
\item Select a candidate using the rule below, assign its team, and repeat.
\end{enumerate}
\begin{center}
\begin{tabularx}{\linewidth}{@{}lX@{}}
\toprule
Policy & Selection rule \\
\midrule
Nearest-first & Smallest distance from the candidate target to its closest free drone \\
Type-priority & Type 1 before Type 2 before Type 3; break ties by the same distance \\
Random & Independently sample an available target uniformly; retain ongoing strikes \\
\bottomrule
\end{tabularx}
\end{center}
A candidate team requires enough free drones to fill the unreserved slots.

\lead{Score resolver}
When several drones choose the same target, the resolver uses their actor
scores $\ell_t^{ik}$ to allocate the available slots:
\begin{enumerate}
\item Retain drones already participating in a strike.
\item For each contested target $k$, rank the other candidates by
decreasing $\ell_t^{ik}$ and admit the highest-scoring drones into the remaining slots.
\item Rejected drones mask that target and choose again. Repeat until the
assignment is resolved.
\end{enumerate}
Training collects direct policy choices; evaluation applies the score resolver.

\newpage
\subsection{Learning curves on held-out scenarios}
\begin{figure}[htbp]
\centering
\IfFileExists{runs/target_coma_v15_train_dist_audit/learning_curves_iqm_ma5.pdf}{%
  \includegraphics[width=\linewidth]{runs/target_coma_v15_train_dist_audit/learning_curves_iqm_ma5.pdf}%
}{%
  \fbox{\parbox[c][5.5cm][c]{0.93\linewidth}{\centering
  Learning-curve figure\\[6pt]
  Left: episode return IQM\\
  Right: success rate relative to the oracle}}%
}
\caption{Learning curves on 100 fixed held-out scenarios sampled from the
training distribution. The two learned-policy curves use the same checkpoints
with score-based or distance-based conflict ranking. Returns are summarized
by the interquartile mean (IQM); success rates are divided by the oracle's
success rate. Learned-policy curves use a trailing five-checkpoint average.}
\label{fig:learning-curves}
\end{figure}

\lead{Evaluation statistics}
\begin{center}
\begin{tabularx}{\linewidth}{@{}lX@{}}
\toprule
Element & Meaning \\
\midrule
Horizontal axis & Training transitions accumulated over active drones \\
Return IQM & Mean of the middle 50\% of the 100 episode returns \\
Relative success & Policy success rate divided by the oracle success rate of 58\% \\
Score / distance ranking & Two conflict-resolution rules applied to the same trained actor \\
Nearest-first / type-priority & Rule-based target-assignment baselines \\
Oracle & Physical upper bound on attainable target value \\
\bottomrule
\end{tabularx}
\end{center}


\begin{itemize}
\item The score-ranked policy improves rapidly
during the first few million transitions, then remains near a return IQM of
9.2 and approximately 90\% of the oracle success rate.
\item The score-ranked and distance-ranked
curves become closely aligned. For competing drone--target assignments with
the same target value, this is consistent with learned scores favoring
shorter distances. Distance ranking provides a comparison of how the same
actor's competing choices are resolved.
\item At 10 million transitions, the plotted
score-ranked policy reaches a return IQM of 9.198 and relative success of
89.0\%. Type-priority reaches 9.12 and 87.9\%; nearest-first reaches 8.88
and 75.9\%. The oracle return IQM is 9.84.
\end{itemize}

\section{Scenario notes}
Three configurations to keep in mind when examining assignments:
\begin{center}
\begin{tabularx}{\linewidth}{@{}cXlc@{}}
\toprule
Case & Location of both Type 1 targets & First centre distance & $B$ \\
\midrule
1 & Second formation & 3.5--4.0 km & 5 \\
2 & Nearby first formation & 2.5--2.7 km & 12 \\
3 & Distant first formation & 5.4--5.8 km & 12 \\
\bottomrule
\end{tabularx}
\end{center}
Distances are measured from the friendly formation centre.

\end{document}
